% Created 2025-09-26 Fri 22:33
% Intended LaTeX compiler: pdflatex
\documentclass[11pt]{article}
\usepackage[utf8]{inputenc}
\usepackage[T1]{fontenc}
\usepackage{graphicx}
\usepackage{longtable}
\usepackage{wrapfig}
\usepackage{rotating}
\usepackage[normalem]{ulem}
\usepackage{amsmath}
\usepackage{amssymb}
\usepackage{capt-of}
\usepackage{hyperref}
\usepackage[margin=1in]{geometry}
\usepackage{microtype}
\usepackage{graphicx}
\usepackage{booktabs}
\usepackage[sorting=none,backend=biber]{biblatex}
\addbibresource{./refs.bib}
\author{Matthew Younger - Pledged}
\date{\today}
\title{From Sound to Story\\\medskip
\large Exploratory Data Analysis of Morse Code Signals in R}
\hypersetup{
 pdfauthor={Matthew Younger - Pledged},
 pdftitle={From Sound to Story},
 pdfkeywords={},
 pdfsubject={},
 pdfcreator={Emacs 30.1 (Org mode 9.7.11)}, 
 pdflang={English}}
\begin{document}

\maketitle
\setcounter{tocdepth}{2}
\tableofcontents

\section{Summary}
\label{sec:orgf0ac273}
The purpose of this project is to satisfy the requirements for a College-level
Data Science project which seeks to explore the challenges of decoding CW Morse Code, one of the oldest
forms of digital communication.
\section{Introduction and Background}
\label{sec:org21c3053}
Morse Code is transmitted as short and long tones \emph{(dits and dahs)} separated by precisely-timed periods of silence.
While decoding Morse Code automatically in noisy recordings is a challenging problem, this project focuses on a simpler, more foundational task:
applying basic data science techniques in R to explore and visualize Morse Code audio. Rather than training models or implementing complex
algorithms, the goal is to explore novel techniques of cleaning, transforming, and interpreting raw audio data.
\section{Planned Scope}
\label{sec:org85b2d44}
\subsection{Goals}
\label{sec:orgb48aa8d}
The primary focus of the project, at least for the duration of the course, will be on conducting EDA and data visualization. As \textbf{stretch goals,}
It would be beneficial to the \texttt{AudioCW project} \autocite{audioCW} as a whole to continue exploring Machine Learning and using the data obtained from analyzing the data for
training the model. As such, the primary goals are as follow:

\begin{enumerate}
\item Apply basic pre-processing to Morse Code audio recordings (denoising, filtering)
\item Generate spectrograms and visualize audio patterns
\item Compute acoustic indices such as Acoustic Complexity Index (ACI) and Acoustic Richness (AR)
\item Extract simple features, such as average tone duration, for estimating the Words Per Minute of the transmission.
\item Create clear visualizations of the underlying data of interest in each file.
\item Generate a dataset that would be the basis for a future Machine Learning project.
\end{enumerate}
\subsection{Constraints}
\label{sec:orgc4ccdf6}
As of present, this is an entirely new area of research for the author. Additionally, they are the sole researcher for this project.
As such, there are many factors that may negatively influence the outcome of the study arising from inexperience, ignorance, or simply from
exploring a new subject in a relative vacuum. It's entirely possible, or perhaps rather likely that there will be misconceptions along the way
and the result could be diminished by this being a solo project.
\section{Methodology}
\label{sec:org7ed48ba}
\subsection{Visualization and EDA}
\label{sec:org3b98012}
The proposed approach is to use computer-generated training data originally intended for a \emph{Morse Learning Machine Challenge} from Kaggle \autocite{ag1le2014},
which will be analyzed in R using various candidate libraries to conduct \textbf{Exploratory Data Analysis} using the following R libraries:
\begin{itemize}
\item Explore the potential to apply bioacoustic analysis more broadly to digital signal processing using \texttt{Sonic Screwdriver} \autocite{ssdriverCran}
\item Calculate the Acoustic Indices using \texttt{soundecology} \autocite{soundecology}
\item Conduct basic sound analysis with \texttt{Seewave}  \autocite{seewaveCran}
\item Utilize \texttt{R} Programming Language \autocite{RLang}, \texttt{GNU Emacs} \autocite{emacs}, and \texttt{GNU Org Mode} \autocite{orgmode}
as the main programming development environment.
\end{itemize}
\subsubsection{Proposed Visualization Techniques:}
\label{sec:org5e8ddda}
\begin{enumerate}
\item Full-color spectrograms and waveforms of relative amplitude envelope, annotated tone on/off segments.
\item Timeline of gaps per element, char, and word.
\item WPM estimate computed and annotated per the PARIS standard.
\item Noise/SNR estimate for each recording.
This information would be very valuable, as a high noise floor means any algorithm or
conventional means of decoding the signal will presumably fail under sufficiently unfavorable conditions.
\end{enumerate}
\subsection{Machine Learning}
\label{sec:org4770202}
In the unlikely case that there is time permitting to explore the \textbf{flex goals,} this research could extend into Machine Learning and interact with other software in the parent project.
Conceptually, there are several machine learning techniques that could be applied:

\begin{itemize}
\item \textbf{Naive Bayes Classification:}
The author of the Kaggle dataset previously attempted to use Bayesian probability to estimate the \emph{probability} that
a given tone is a dit/da or even an entire word fits inside a given \emph{dictionary} of words, \autocite{a1gleblog}. This could potentially
be used by itself, or as a type of post-processing after a simpler and lighter algorithm does the bulk of the
heavy lifting. Many practical conversations are very formulaic and predictable, due to the nature of HAM radio;
The typical conversation usually consists of a callsign, a brief exchange of location, reception clarity, contest details,
type of equipment being used to transmit and receive with, and an end of transmission "goodbye." The small sample space makes
this a good option - even more so in the case of a robot that's exchanging a highly formulaic message sequence.

\item \textbf{Clustering:}
There's promising literature that indicates clustering is very good for identifying bird sounds in audio recordings,
particularly when combined with uniform manifold approximation and dimension reduction, \autocite{RSoundAnalysis}.
Just as birds have a distinct sound, so does each letter in CW. Per the guidance of the largest CW enthusiasts club in the country,
humans learn best by recognizing entire sounds, as opposed to decoding the sequence of dashes and dits on the fly.

Literature from the Long Island CW enthusiasts club, one of the largest clubs dedicated to morse code in the United States,
emphasizes the effectiveness of having students identify entire acoustic sounds, rather than \emph{counting} the sequenes, \autocite{LICWTrainingMaterials}.
The group derives this understanding from careful review of Psychologist Ludwig Koch's study of learning Morse in 1936, having
reached this conclusion after having an external entity translate Koch's work from the original German \autocite{KochReport}.
\end{itemize}
\section{Timeline}
\label{sec:orgb25297c}
Per the course outline, this project is to be managed under the \textbf{scrum project management} framework,
where we are encouraged as aspiring Data Scientists to continually \emph{evaluate} and \emph{learn through our experiences}
with attention to our personal \textbf{skill stack and problem-solving abilities} as they relate to the task at hand.

Particularly, the project will be divided into four checkpoints:

\begin{enumerate}
\item \textbf{This proposal stage:} where the idea is pitched to a client or end-user.
\item \textbf{Literature Review}: wherein we evaluate other scholarly material to gain helpful insights.
A core skill in good Science is having an adequate grasp of what others have tried before.
\item \textbf{Abstract:} Likewise, it is important that one's research is clearly articulated. An abstract helps
other readers to efficiently determine whether our work is relevant to their studies.
\item \textbf{Presentation:} Finally, the conclusion of our study is presented to other scientists, students, colleagues,
or more universally an end-user. Almost always done in a live "show-and-tell" style, this an opportunity to showcase
a project and generate further interest to continue development or gain patronage for a future project (or to pass a class, in this case).
\end{enumerate}
\section{Desired Outcomes}
\label{sec:org47efe15}
To aid in the visual evaluation, openly-available audio software, such as Tenacity \autocite{tenacityFOSS}, can be used to visually and audibly compare our generated
visuals to that of a commonly-used audio tool. Synthetically-generated data can be interpreted visually in this way when it's pure tones without any noise or distortion,
which so too should this data if it's adequately cleaned and coherently represented. Given that producing clean visual graphs and representations is in-step with
cleaning the audio itself, it is possible then to listen to the cleaned test sets and compare the provided labels given in the project manually. Admittedly, this is a somewhat
manual process that at scale would be infeasible. However, at present under the present conditions, a better alternative is not apparent. Similarly, a clean waveform can be visually
examined and the code extracted by a similar manual means. If the reviewer is able to understand the code once it's been cleaned, it could be argued that the goal has been fulfilled,
although this also is rather empirical.

Should this project advance to the stage of providing a machine learning algorithm capable of decoding the audio, the candidate Kaggle dataset includes a script that measures the
\emph{Levenshtein distance}, by which the challenge author claims the results can be objectively evaluated.

Overall evaluation will be conducted formally by the course instructor per the University's grading policy,
the course syllabus, and the grading rubric as provided for each stage. In addition, less informal but no less
important evaluation is conducted by the researcher, per the aforementioned agile framework.
\section{Deliverable}
\label{sec:orgb044501}
The deliverable of this project will ideally be a cleaned dataset and a clearly-defined methodology for overcoming the difficulties that the
noise and speed of keying in this dataset present. This will be hosted in a Git repository and will be publicly available for others to review.
Additionally, the structure of the semester's exercise should yield a coherent presentation and abstract that details the results of the
exercise as a whole, which may guide others in similar pursuits.

The literature that this author has found so far seems to indicate that audio processing methods from the novel field of ecoacoustics can perhaps
be adapted to better understand other types of audio. Given that simplicity and affordability of collecting audio and its abundance,
any research that successfully furthers the pursuit of understanding audio with data science techniques should be applauded.
\section{Conclusion}
\label{sec:org36c5a97}
Foremost, it is this project's intention to examine CW/Morse Code audio emissions gathered by microphone via digital signal processing
from a Data Science perspective. Audio files themselves are data structures which contain meaningful data, as well as outlier data: static, or literal noise, which can be
\textbf{cleaned, manipulated, analyzed, and visualized} in order to \textbf{extract meaningful insights and information:} an encoded message.

Ultimately, Data Science in this context will aid the parent program's ability to decode raw data from a microphone into actionable information that can either be displayed to an end-user
and/or executed directly by another program in an application stack as part of a larger project. As such, if this project can meet its goals, it not only adequately demonstrates
Data Science's intrinsic value in analyzing and understanding complicated information such as audio waveforms, but additionally showcases its value as an applied science to further other
other fields, such as Robotics, Telecommunications, and Machine Learning, which are likewise key areas relating to the larger project.
\section{References}
\label{sec:orge19f71a}
\printbibliography
\end{document}
